\documentclass[../../../include/open-logic-section]{subfiles}

\begin{document}

\olfileid{sth}{choice}{banach}
\olsection{বানাখ--তার্স্কি কূটাভাস}

বুলস যেমন প্রতিস্থাপনের ন্যায্যতা দেখাতে
চেয়েছিলেন, তেমন চয়নের ন্যায্যতাও
\emph{বাহ্যিক} বিবেচনায় দেখার চেষ্টা করা যায়
(দেখো \olref[replacement][extrinsic]{sec})।
চয়ন নিলে অনেক কাম্য ফল পাওয়া যায়:
সব অঙ্কবাচক সংখ্যার তুলনা, প্রত্যেক সেটের
সুক্রমবিন্যাস, বিশেষ ধরনের ক্রমসংখ্যা হিসেবে
অঙ্কবাচক সংখ্যাকে ধরা, ইত্যাদি।

তবে কখনো বলা হয়, চয়নের \emph{অকাম্য} ফলও
আছে। প্রধানত \cite{BanachTarski1924}-এর
একটি ফলের জন্য এ কথা বলা হয়।

\begin{thm}[বানাখ--তার্স্কি কূটাভাস ($\ZFC$-এ)]
যেকোনো বলকে সসীমসংখ্যক খণ্ডে ভাগ করে,
সেগুলি ঘুরিয়ে ও স্থানান্তর করে আবার সাজালে
সেই বলের দুটি প্রতিলিপি গঠন করা যায়।
\end{thm}
\noindent
প্রথম দেখায় এটি বিস্ময়কর। দুটি বলের আয়তন
মূল বলটির \emph{দ্বিগুণ}। কিন্তু অনমনীয় গতি—
ঘোরানো ও স্থানান্তর—আয়তন বদলায় না। তাই
মনে হয়, বানাখ--তার্স্কি বুঝি শূন্য থেকে নতুন
পদার্থ সৃষ্টি করতে দেয়।

আরও আশ্চর্য ফল আছে।\footnote{দেখো
\citet[Theorem 3.12]{Wagon2016}।}
একই ধরনের যুক্তিতে একটি মটরশুঁটি সসীমসংখ্যক
খণ্ডে কেটে সেগুলি ঘুরিয়ে ও স্থানান্তর করে
বিগ বেনের আকার ও আয়তনের বস্তু বানানো যায়।

তবে কেবল $\ZF$ থেকে এগুলির কোনোটিই
অনুসৃত হয় না।\footnote{যদিও চয়নের চেয়ে
কঠোরভাবে দুর্বল কিছু নীতিতে বানাখ--তার্স্কি
প্রমাণ করা যায়; দেখো \citet[303]{Wagon2016}।}
তাই সিদ্ধান্তের মুখে পড়ি: চয়ন প্রত্যাখ্যান করব,
নাকি ``কূটাভাস'' মেনে নিতে শিখব?

আমাদের পরামর্শ, ``কূটাভাস''টি মেনে নিতে
শেখা উচিত। বস্তুত আমাদের কাছে এটি খুব একটা
কূটাভাসই মনে হয় না। নিচের ফলগুলির তুলনায়
একে বেশি বা কম কূটাভাসমূলক ভাবার কারণ
দেখি না:\footnote{\citet[276--7]{Potter2004},
\citet[16]{Weston2003}, এবং
\citet[31, 308--9]{Wagon2016} অন্য
উদাহরণে কাছাকাছি কথা বলেছেন।}
\begin{enumerate}
	\item $(0,1)$ অন্তরালের বিন্দু যত, $\Real$-এও তত।
	% BN-SRC-471: remove the unmatched closing parenthesis in the tangent map.
	\\\emph{প্রমাণ}: $\tan(\pi(r-\nicefrac{1}{2}))$ দেখো।
	\item সরলরেখায় যত বিন্দু, বর্গক্ষেত্রেও তত।
	\\দেখো \olref[his][set][pathology]{sec} এবং
	\olref[his][set][cantorplane]{sec}।
	\item স্থান-পূরণকারী বক্ররেখা আছে।
	\\দেখো \olref[his][set][pathology]{sec} এবং
	\olref[his][set][hilbertcurve]{sec}।
\end{enumerate}
এই তিন ফলের কোনোটি প্রমাণে চয়ন লাগে না।
আজ এগুলিকে গণিতের বিস্ময়কর, সুন্দর ফল
বলি। হয়তো বানাখ--তার্স্কি কূটাভাসের
প্রতিও একই মনোভাব নেওয়া উচিত।

অবশ্য একটি কারিগরি বিষয় বুঝতে হবে;
শান্তভাবে ভাবলেই হয়। অনমনীয় গতি আয়তন
অক্ষত রাখে। তাই বলের পাঁচটি\footnote{আমরা
কূটাভাসটি ``সসীমসংখ্যক খণ্ড'' দিয়ে বলেছি।
বস্তুত \citet{Robinson1947} দেখিয়েছিলেন,
\emph{পাঁচটি} খণ্ডেই বিভাজন সম্ভব, কমে নয়।
প্রমাণ দেখো \citet[pp.~66--7]{Wagon2016}।}
খণ্ডের সবগুলি \emph{পরিমাপযোগ্য} হতে পারে না।
মোটামুটি, খণ্ডগুলির প্রত্যেকটির আলাদা আয়তন
দেওয়ার অর্থ নেই। এগুলিকে বিন্দুর এমন
``অসীমভাবে ছড়ানো'' সমাবেশ ভাবা যায়,
যার ছবি আঁকা যায় না। এই সেটগুলির অস্তিত্ব
``অদ্ভুত'' মনে হতে পারে। কিন্তু \stagesacc{}-এ
প্রকাশিত আগে থেকে পাওয়া সেটগুলির \emph{সব
সম্ভাব্য} সংকলন গঠনের নির্দেশ থেকেই যেন
এদের অস্তিত্ব বেরিয়ে আসে।

তাতেও রাজি না হলে বানাখ--তার্স্কি গ্রহণের
পক্ষে একটি শেষ (বাহ্যিক) যুক্তি আছে। এটি
গণিতের সেরা রসিকতাটি সঙ্গে সঙ্গে দেয়:
\begin{enumerate}
	\item[] \emph{প্রশ্ন}। ``Banach--Tarski''-র
	অক্ষরগুলি নতুন করে সাজালে কী হয়?
	\item[] \emph{উত্তর}। ``Banach--Tarski Banach--Tarski''।
\end{enumerate}

\end{document}
